% Unit 8 map -- one page.
\documentclass{apchem-worksheet}

\apcunit{8}
\apcunittitle{Acids and Bases}
\apcdoctitle{Unit Map}
\apcsubtitle{CED Unit 8 \textbullet\ \textbf{11--15\%} of the AP Exam \textbullet\ 9 blocks + float + exam}

\begin{document}
\apctitleblock

\begin{apcnote}[How this unit is graded --- and what makes it different]
Unit exam \textbf{Fri Feb 19}: 30 multiple choice (\textbf{no calculator})
$+$ 1 long and 1 short free-response, 76 minutes, 44 points.

At \textbf{11--15\%} this is the second-heaviest unit in the course, behind
only Unit 3. It is also the longest --- nine blocks --- because it is really
three units stacked: weak-acid equilibria, buffers, and titrations.

None of it is new machinery. A weak acid problem is a Unit 7 ICE table with
$K_a$ instead of $K$. A buffer is that same equilibrium with the conjugate
base already present. A titration is a \emph{stoichiometry} problem followed
by whichever equilibrium survives. The skill being tested is deciding
\textbf{which of those four situations you are in} before you calculate
anything.
\end{apcnote}

\section*{\sffamily Block calendar}

\renewcommand{\arraystretch}{1.25}
\noindent\begin{tabular}{@{}p{2.1cm}p{4.9cm}p{1.6cm}p{3.9cm}p{1.9cm}@{}}
\toprule
\textbf{Date} & \textbf{Topics} & \textbf{CED} & \textbf{Read (PDF pp.)} & \textbf{Practice}\\
\midrule
Wed Jan 27 & Br\o nsted--Lowry; conjugate pairs; $K_a$ & 8.1 & \S14.1--14.2 \textbullet\ 697--705 & WS 1\\
Fri Jan 29 & pH and pOH; strong acids and bases & 8.2 & \S14.3--14.4, 14.6 \textbullet\ 705--724 & WS 1\\
Mon Feb 1 & Weak acids: $K_a$ and ICE & 8.3 & \S14.5 \textbullet\ 710--719 & WS 2\\
Wed Feb 3 & Weak bases; polyprotic acids & 8.3 & \S14.6--14.7 \textbullet\ 719--730 & WS 2\\
Fri Feb 5 & Structure; p$K_a$; salt solutions & 8.6, 8.7 & \S14.8--14.9 \textbullet\ 730--740 & WS 3\\
Mon Feb 8 & Acid--base reactions; what a buffer is & 8.4 & \S15.1--15.2 \textbullet\ 756--765 & WS 3\\
Wed Feb 10 & Buffers and Henderson--Hasselbalch & 8.8, 8.9 & \S15.2 \textbullet\ 763--767 & WS 4\\
Fri Feb 12 & Buffer action and capacity & 8.10 & \S15.2--15.3 \textbullet\ 761--770 & WS 4\\
Mon Feb 15 & Titrations, indicators, pH and solubility & 8.5, 8.11 & \S15.4--15.5 \textbullet\ 770--792; \S16.1 \textbullet\ 812 & WS 5 (FRQ)\\
Wed Feb 17 & Reteach / float & --- & review sheet & ---\\
\textbf{Fri Feb 19} & \textbf{UNIT 8 EXAM} & all & review sheet & ---\\
\bottomrule
\end{tabular}

{\sffamily\footnotesize Dates from the co-op calendar. Unit 9 begins
Mon Feb 22.}

\vspace{0.4em}
\section*{\sffamily By the end of this unit you can\ldots}

\begin{itemize}[leftmargin=*,itemsep=0.1em]
  \item Identify Br\o nsted--Lowry conjugate acid--base pairs and relate
        $K_a$ to $K_b$. \ced{8.1}
  \item Calculate pH and pOH for strong acids and bases. \ced{8.2}
  \item Solve weak acid and weak base equilibria with an ICE table.
        \ced{8.3}
  \item Predict the products of an acid--base reaction and recognize a
        buffer. \ced{8.4}
  \item Analyse a titration curve region by region and choose an indicator.
        \ced{8.5}
  \item Explain acid strength from molecular structure. \ced{8.6}
  \item Use p$K_a$ to rank strength and predict salt behaviour. \ced{8.7}
  \item Explain how a buffer resists pH change. \ced{8.8}
  \item Apply the Henderson--Hasselbalch equation. \ced{8.9}
  \item Compare buffering capacity and design a buffer for a target pH.
        \ced{8.10}
  \item Predict qualitatively how pH affects the solubility of a salt.
        \ced{8.11}
\end{itemize}

\begin{apctrap}[Where students lose points in this unit]
Reporting $-\log x$ as the pH for a weak \emph{base} (it is the pOH);
forgetting to double $[\ce{OH-}]$ for \ce{Ba(OH)2} and \ce{Ca(OH)2};
skipping the 5\% check; putting added strong base into an ICE table instead
of reacting it to completion first; inverting the
$\log([\text{base}]/[\text{acid}])$ ratio; assuming a titration's
equivalence point is at pH 7; and \emph{calculating} solubility from pH,
which topic 8.11 explicitly excludes.
\end{apctrap}

\end{document}
